% Generated by `taskgen course examples` from dlp.pohlig_hellman (seed 16).
% Do not edit: change examples.yaml and regenerate.
\begin{example}[алгоритм Полига--Хеллмана]\label{ex:ph-full}
Решить задачу дискретного логарифмирования с помощью алгоритма Полига-Хеллмана: \[2^x \equiv 59\ (\textrm{mod}\ 61).\] Решения частных задач и мультипликативно обратные можно находить перебором («на глазок»). Систему необходимо выписать и решить в явном виде.

\par\noindent\textit{Решение.} $\ord(2) = 60$: $2^{60} \equiv 1$, $2^{30} \equiv 60$, $2^{20} \equiv 47$, $2^{12} \equiv 9$. $60 = 2^{2} \cdot 3 \cdot 5$. Для каждого $q^e \parallel 60$: $g_i = g^{60/q^e}$, $h_i = h^{60/q^e}$, решаем $g_i^{x_i} = h_i$ перебором, тогда $x \equiv x_i\ (\textrm{mod}\ q^e)$.\par\noindent $q^e = 4$: $g_i = 2^{15} \equiv 11$, $h_i = 59^{15} \equiv 50$; степени $g_i^k$, $k = 0, 1, \dots$: $1, 11, 60, 50$ $\Rightarrow x_i = 3$.\par\noindent $q^e = 3$: $g_i = 2^{20} \equiv 47$, $h_i = 59^{20} \equiv 47$; степени $g_i^k$, $k = 0, 1, \dots$: $1, 47$ $\Rightarrow x_i = 1$.\par\noindent $q^e = 5$: $g_i = 2^{12} \equiv 9$, $h_i = 59^{12} \equiv 9$; степени $g_i^k$, $k = 0, 1, \dots$: $1, 9$ $\Rightarrow x_i = 1$.\par Система и её решение по КТО:\[\left\{\begin{array}{l}
x \equiv 3\ (\textrm{mod}\ 4) \\
x \equiv 1\ (\textrm{mod}\ 3) \\
x \equiv 1\ (\textrm{mod}\ 5) \\
\end{array}\right.\qquad M = 60.\]
$M_{1} = M/4 = 15 \equiv 3$, $y_{1} = M_{1}^{-1} \equiv 3\ (\textrm{mod}\ 4)$; $M_{2} = M/3 = 20 \equiv 2$, $y_{2} = M_{2}^{-1} \equiv 2\ (\textrm{mod}\ 3)$; $M_{3} = M/5 = 12 \equiv 2$, $y_{3} = M_{3}^{-1} \equiv 3\ (\textrm{mod}\ 5)$.\[\begin{aligned}x &\equiv \textstyle\sum_i r_i M_i y_i \\ &= 3\cdot15\cdot3 + 1\cdot20\cdot2 + 1\cdot12\cdot3 \\ &\equiv 31\ (\textrm{mod}\ 60).\end{aligned}\]
\[x = 31\ (\textrm{mod}\ 60),\quad 60 = 2^{2} \cdot 3 \cdot 5.\]
\end{example}
